\begin{abstract}
\noindent
Chất lượng dữ liệu ảnh hưởng trực tiếp đến hiệu năng của mô hình học máy, nhưng mức ảnh hưởng thay đổi theo tác vụ và thuật toán. Các công cụ hiện có, tiêu biểu là DQSOps, gộp nhiều chiều chất lượng thành một điểm duy nhất, giống nhau cho mọi tác vụ và không gắn với hiệu năng thật của mô hình phía sau. Đề tài xây dựng một khung đánh giá theo hướng ngược lại: đo chất lượng dữ liệu bằng một profile năm chiều (Completeness, Feature Accuracy, Target Accuracy, Uniqueness, Class Balance), rồi học một meta-model Gradient Boosting dự đoán \emph{mức suy giảm hiệu năng thật} ($\Delta$F1, $\Delta$R$^2$, $\Delta$AMI) của mô hình downstream, có biết loại tác vụ. Dữ liệu huấn luyện meta-model gồm 312 lượt thí nghiệm làm bẩn có kiểm soát trên ba bộ dữ liệu (HIGGS, Beijing Air Quality, Covertype) với sáu thuật toán. Với đánh giá chéo theo nhóm để chống rò rỉ, meta-model đạt MAE 0,055, thấp hơn weighted score (0,077) và baseline kiểu DQSOps (0,084). Mức chênh so với DQSOps có ý nghĩa thống kê theo bootstrap; so với weighted score thì chưa. Meta-model cũng xếp hạng mức suy giảm tốt nhất (Spearman 0,66 so với 0,32 và 0,06). Phân tích SHAP cho thấy thứ tự quan trọng của các chiều chất lượng khác nhau theo tác vụ, trái với bộ trọng số cố định của PCA. Về quy mô, bản profiling bằng PySpark cho cùng kết quả với pandas đến 11 triệu dòng và vẫn chạy được ở 22 triệu dòng, mốc mà pandas hết bộ nhớ.
\end{abstract}
